\section{Preliminaries}

We will introduce notation, 
and then we will discuss two relevant works---the papers by 
Kalhan, Makarychev, and Zhou \citeyearpar{KMZ19} and Davies, Moseley, and Newman \citeyearpar{DMN23}.

\subsection{Notation} \label{sec: notation}
Recall our input to the correlation clustering problem is 
$G=(V,E)$, an unweighted, complete graph on $n$ vertices, 
and every edge is assigned a label of either positive $(+)$ or negative $(-)$.
Let the set of positive edges be denoted $E^+$ and the set of negative edges $E^-$.
Then, we can define the \emph{positive neighborhood} and \emph{negative neighborhood} of a vertex $u$ as $N_u^+ =\{v\in V \mid (u,v) \in E^+\}$ and $N_u^- =\{v\in V \mid (u,v) \in E^-\}$, respectively.
We further assume without loss of generality that every vertex has a positive self-loop to itself.

A \emph{clustering} $\mathcal{C}$ is a partition of $V$ into \emph{clusters} $C_1,\ldots,C_k$ (but recall that $k$ is \textit{not} pre-specified). 
Let $C(u)$ denote the cluster that vertex $u$ is in, i.e., if $\mathcal{C}$ has $k$ clusters, there exists exactly one $i \in [k]$ such that $C(u) = C_i$.
It is also helpful to consider the vertices in a different cluster than $u$, and so we let $\overline{C(u)} = V \setminus C(u)$ denote this.
We say that a positive edge $e=(u,v) \in E^+$ is a \emph{disagreement} with respect to $\mathcal{C}$ if $v \in \overline{C(u)}$. On the other hand, we say that a negative edge $e=(u,v) \in E^-$ is a disagreement with respect to $\mathcal{C}$ if $v \in C(u)$.
For a fixed clustering $\mathcal{C}$, we denote the \emph{disagreement vector} of $\mathcal{C}$ as $y_{\mathcal{C}} \in \mathbb{Z}_{\geq 0}^n$, where for $u \in V$, $y_{\mathcal{C}}(u)$ is the number of edges incident to $u$ that are disagreements with respect to $\mathcal{C}$. We omit the subscript throughout the proofs when a clustering is clear.

Throughout, we let $\textsf{OPT}$ be the optimal objective value, 
and the $\ell_p$-norm to which it corresponds will be clear from context.
The next fact follows from the definitions seen so far (recalling also the positive self-loops). 



\begin{fact} \label{pairwise_intersections}
For any $u,v \in V$, possibly with $u=v$, $n = |N_u^+ \cap N_v^+| + |N_u^- \cap N_v^-| + |N_u^+ \cap N_v^-| + |N_u^- \cap N_v^+|.$
\end{fact}





\subsection{Summary of work by Kalhan, Makarychev, and Zhou}
The standard linear program relaxation for correlation clustering 
is stated in LP \ref{KMZ_LP}.\footnote{Technically this is a convex program as the objective is convex. For simplicity we will refer to it as an LP as the constraints are linear.}
In the \textit{integer} LP, the variable $x_{uv}$ indicates whether vertices $u$ and $v$ will be in the same cluster (0 for yes, 1 for no), and the disagreement vector is $y$; the optimal solution to the integer LP has value \textsf{OPT}, while the optimal solution to the relaxation gives a lower bound on $\textsf{OPT}$. 
Note the triangle inequality is enforced on all triples of vertices, inducing a semi-metric space on $V$.
Throughout this paper, as in \cite{DMN23}, we refer to the algorithm by Kalhan, Makarychev, and Zhou as the KMZ algorithm. 
The \emph{KMZ algorithm} has two phases: it solves LP \ref{KMZ_LP}, 
and then uses the \textit{KMZ rounding algorithm} (Algorithm \ref{KMZ-alg}) to obtain an integral assignment of vertices to clusters. 
See Appendix \ref{sec: lp_rounding_alg} for its formal statement.
At a high-level, the KMZ rounding algorithm is an iterative, ball-growing algorithm
that uses the semi-metric to guide its choices on forming clusters.
Their algorithm is a 5-approximation, and produces different clusterings for different $p$, since the optimal solution $x^*$ to LP \ref{KMZ_LP} depends on $p$.



\vspace{-1cm}
\begin{flushleft}
\begin{equation*}\label{KMZ_LP}
\end{equation*}
LP \ref{KMZ_LP}
\end{flushleft}
\vspace{-1cm}
\begin{align}
     &\min {\lVert y \rVert}_p \notag\\
    \textsf{s.t. } y_u &= \sum_{v \in N_u^+} x_{uv} + \sum_{v \in N_u^-} (1-x_{uv})  && \forall u \in V \notag\\
     x_{uv}&\leq x_{uw}+x_{vw}  &&\forall u,v,w \in V \notag\\
     0 &\leq x_{uv} \leq 1  &&\forall u,v \in V. \notag
\end{align}

\begin{definition} \label{def: frac_cost}
    Let $f$ be a semi-metric on $V$, i.e., taking $x=f$ gives a feasible solution to LP \ref{KMZ_LP}. The \emph{fractional cost of $f$ in the $\ell_p$-norm objective} is the value of LP \ref{KMZ_LP} that results from setting $x=f$. When $p$ is clear from context, we will simply call this the fractional cost of $f$. 
\end{definition}



\subsection{Summary of work by Davies, Moseley, and Newman}\label{sec: DMN-summary}


The main take-away from the work of Kalhan, Makarychev, and Zhou \citeyearpar{KMZ19} is that one only requires a semi-metric on the set of vertices,
whose cost is comparable to the cost of an optimal solution,
as input to the KMZ rounding algorithm (Algorithm \ref{KMZ-alg}). 
Thus, the insight of Davies, Moseley, and Newman \citeyearpar{DMN23} for the $\ell_{\infty}$-norm objective is that one can combinatorially construct such a semi-metric without solving an LP, and at small loss in the quality of the fractional solution.
They do this by introducing the \emph{correlation metric}.


\begin{definition}[\cite{DMN23}] \label{def: orig_corr_metric}
For all $u,v \in V$, the \textit{\dn}~defines the 
distance between $u$ and $v$ as
\begin{align}
 \label{eq: equiv_def}
d_{uv} &= 1-\frac{|N_u^+ \cap N_v^+|}{|N_u^+ \cup N_v^+|} \notag \\
&=  \frac{|N_u^+ \cap N_v^-| +|N_u^- \cap N_v^+| }{|N_u^+ \cap N_v^+| + |N_u^+ \cap N_v^-| + |N_u^- \cap N_v^+|}.
\end{align}
\end{definition}

Note that the rewrite in Line (\ref{eq: equiv_def}) is apparent from Fact \ref{pairwise_intersections}.


The correlation metric captures useful information  succinctly.
Intuitively, if $u$ and $v$ have relatively large positive intersection, i.e., $N_u^+ \cap N_v^+$ is large compared to their other relevant joint neighborhoods $\left (N_u^+ \cap N_v^-\right ) \cup \left ( N_u^- \cap N_v^+\right )$, then from the perspective of $u$ and $v$,
fewer disagreements are incurred by putting
$u$ and $v$ in the same cluster than by putting them in different clusters.
This is because if $u$ and $v$ are in the same cluster, then they have disagreements on edges $(u,w)$ and $(v,w)$ for $w \in \left (N_u^+ \cap N_v^- \right )\cup \left (N_u^- \cap N_v^+\right )$, but if they are in different clusters, then $u$ and $v$ have disagreements on edges $(u,w)$ and $(v,w)$ for  $w \in N_u^+ \cap N_v^+$. 
Note that the metric is not normalized by $n$ (the size of the joint neighborhood), but instead by $n-| N_u^- \cap N_v^- | = |N_u^+ \cap N_v^+| + |N_u^+ \cap N_v^-| + |N_u^- \cap N_v^+|$.
In addition to the above intuition on the normalization factor, we also observe that $w \in N_u^- \cap N_v^- $ do not necessarily force disagreements on $(u,w)$ and $(v,w)$, since $w$ can go in a different cluster than both $u$ and $v$ without penalty.
For more on intuition behind the correlation metric, see Section 2 in \cite{DMN23}.

Davies, Moseley, and Newman \citeyearpar{DMN23} prove that the correlation metric $d$ can be used as input to the KMZ rounding algorithm by showing that (1) $d$ satisfies the triangle inequality and (2) the fractional cost of $d$ in the $\ell_{\infty}$-norm (recall Definition \ref{def: frac_cost}) is no more than 8 times the value of the optimal integral solution ($\textsf{OPT}$). Since the KMZ rounding algorithm (Algorithm \ref{KMZ-alg}) loses a factor of at most 5, inputting $d$ to that algorithm returns a 40-approximation algorithm. 
A benefit of the correlation metric is that it can be computed in time $O(n^{\omega})$, and even faster when the subgraph on positive edges is sparse. 




\subsection{Technical overview}
 It is not hard to see that the correlation metric cannot be used as input to the KMZ algorithm for $\ell_p$-norms other than $p = \infty$, as one cannot bound the fractional cost of the correlation metric against the optimal with only an $O(1)$-factor loss. To see why, consider the star again, as in Figure \ref{fig: star}. Here, for all $u,v \in \{v_1,\ldots,v_{n-1}\}$, $d_{uv} = 1- 1/(n-(n-3)) = 2/3$, but for the $\ell_1$-norm, we need the semi-metric to have the value $1-d_{uv}$ be close to 0, i.e. $O(1/n)$, for such $u,v$, in order for the fractional cost to be comparable to the value of $\textsf{OPT}$ for $p=1$. 


There are several possible fixes one could try to make to the correlation metric. 
One idea is that since one can interpret the correlation metric as a coarse approximation of the probability the Pivot algorithm separates $u$ and $v$\footnote{Pivot operates as follows: Choose a random unclustered $u \in V$. Take $u$ and all its unclustered positive neighbors and let this be the newest cluster. Continue until all vertices in $V$ are clustered.}, 
one could try to adapt the correlation metric to more accurately approximate this probability.\footnote{If in LP \ref{KMZ_LP}, $x_{uv}$ is set exactly to the probability that $u$ and $v$ are separated by Pivot, then $x$ will be a feasible solution with cost at most $\textsf{3OPT}$. However, this probability seems difficult to express in closed form, even approximately.}
Another idea, inspired by an observation below, is that one could define a semi-metric for edges in $E^+$ and another semi-metric for edges in $E^-$, but then there is the difficulty of showing the triangle inequality holds when positive and negative edges are mixed together. 
Both of these ideas were, for us, unsuccessful.


Instead, the following two observations of how the correlation metric works with respect to the $\ell_1$-norm led us to an effective adaptation: 
\begin{enumerate}
\item One can bound the fractional cost, \textit{restricted to positive edges}, of the correlation metric in the $\ell_1$-norm  by an $O(1)$-factor times the optimal solution's cost (see Claim \ref{claim: l1-pos-cm} in Appendix \ref{sec:l1-norm}). Negative edges still pose a challenge.
\item If the subgraph of positive edges is regular, then we show one can bound the fractional cost of the correlation metric in the $\ell_1$-norm on negative edges as well as positive.\footnote{In contrast, one cannot bound the fractional cost of the correlation metric on the (irregular) star (Figure \ref{fig: star}).} See Section \ref{sec: regular_warmup}. 
\end{enumerate}


These observations led us to ask whether some adjustments to the correlation metric might yield a semi-metric with bounded fractional cost in the $\ell_1$-norm or even the $\ell_p$-norm more generally (while still remaining bounded in the $\ell_{\infty}$-norm). Moreover, since the KMZ rounding algorithm does not depend on $p$ (whereas in the KMZ algorithm, the solution to the LP \textit{does} depend on $p$), inputting the same semi-metric to the rounding algorithm produces the same clustering for all $\ell_p$-norms! 

We are ready to define the \emph{adjusted correlation metric}. 
Let $\Delta_u$ denote the positive degree of $u$ (the degree of $u$ in the subgraph of positive edges).
\begin{definition} \label{def: adj_corr_metric}
    Define the \emph{adjusted correlation metric} $f: E \rightarrow [0,1]$ as follows:
    \begin{enumerate}
        \item For $d$ the correlation metric, i.e., 
        $d_{uv} = 1 - \frac{|N_u^+ \cap N_v^+|}{|N_u^+ \cup N_v^+|}$, initially set $f = d$. 
        \item If $e \in E^-$ and $d_e > 0.7$, set $f_e = 1$ (round up). 
        \item For $u \in V$ such that $|N_u^- \cap \{v: d_{uv} \leq 0.7\}| \geq \frac{10}{3}\Delta_u $, set $f_{uv}=1$ for all $v \in V \setminus \{u\}$.
    \end{enumerate}
\end{definition}

The idea in Step 3 is that if the fractional cost of negative edges incident to $u$ is sufficiently large, we instead trade this for the cost of positive disagreements, as 
the rounding algorithm will now put $u$ in its own cluster.
For the $\ell_\infty$-norm, this trade-off is innocuous.
For $\ell_p$-norms in general, a refined charging argument is needed to show that post-processing $d$ in this way sufficiently curbs the (too large) fractional cost of $d$.

In Section \ref{sec: regular_warmup}, we start with a warm-up exercise 
and show that if the graph on positive edges is regular, then 
the (original) correlation metric $d$ has $O(1)$-approximate fractional cost.
Note this section is not necessary to understanding the rest of the paper, 
but we include it in the main body because we find the proof here very clever!
The main technical result of the paper is Section \ref{sec: proofs}, where we prove Theorem \ref{thm: main-acm} 
by showing that the adjusted correlation metric can be input to the KMZ rounding algorithm.
Namely, we will first show (quite easily) that the adjusted correlation metric satisfies an approximate triangle inequality.
Then, it remains to upper bound the fractional cost of the adjusted correlation metric against $\textsf{OPT}$.
We tackle this with a combinatorial charging argument.   This argument leverages a somewhat different approach from that used in \cite{DMN23} and is simpler than their proof for only the $\ell_\infty$-norm.
The \textit{constant} approximation factor obtained from inputting the adjusted correlation metric to Algorithm \ref{KMZ-alg} is bounded above (and below) by universal constants for all $p$ (this is the worst case and one can get better constants for each $p$).
While we keep the argument for general $\ell_p$-norms in the main body of the paper,
the interested reader may find a simplified proof for the $\ell_1$-norm objective in Appendix \ref{sec:l1-norm}.



\section{A Special Case: Regular Graphs} \label{sec: regular_warmup}
In general, the original correlation metric $d$ (Definition \ref{def: orig_corr_metric}) does not necessarily have bounded fractional cost for the $\ell_1$-norm objective (or more generally for $\ell_p$-norm objectives). So, we use the adjusted correlation metric $f$ (Definition \ref{def: adj_corr_metric}) as input to the KMZ rounding algorithm (Appendix \ref{sec: lp_rounding_alg}). In this section, we show that \textit{if the subgraph of positive edges is regular}, then the correlation metric $d$ can be used as is (i.e., without the adjustments in Steps 2 and 3 of Definition \ref{def: adj_corr_metric}) to yield a clustering that is constant approximate for the $\ell_1$-norm and $\ell_{\infty}$-norm simultaneously:

\begin{theorem}
    Let $G=(V,E)$ be an instance of unweighted, complete correlation clustering, and let $E^+$ denote the set of positive edges. Suppose that the subgraph induced by $E^+$ is regular. Then the fractional cost of $d$ in the $\ell_1$-norm objective is within a constant factor of $\textsf{OPT}$:

    \[\sum_{u \in V} \sum_{v \in N_u^+} d_{uv} + \sum_{u \in V} \sum_{v \in N_u^-} (1-d_{uv}) = O(\textsf{OPT}).\]
    Therefore, the clustering produced by inputting $d$ to the KMZ rounding algorithm is a constant-factor approximation simultaneously for the $\ell_1$-norm and $\ell_{\infty}$-norm objectives. 
\end{theorem}

While this proof is quite different than the one for arbitrary-degree graphs,
we include it because we find it beautiful,
and it gives useful intuition for where the correlation metric breaks down for irregular graphs.

\begin{proof}
    Let $\Delta$ be the (common) degree of the positive subgraph. To show that the fractional cost of $d$ in the $\ell_1$-norm objective is $O(\textsf{OPT})$ for regular graphs, we will use a dual fitting argument. The LP relaxation we consider is from \cite{ACN-pivot}, which uses a dual fitting argument to show constant approximation guarantees for Pivot (although the proof here does not otherwise resemble the proof for Pivot). The primal is given by 

\begin{equation} \label{eq: lp_pivot} \tag{$P$}
 \min\Big\{\sum_{e \in E} x_e \mid x_{ij}+x_{jk}+x_{ki} \geq 1, \forall ijk \in \mathcal{T}, x \geq 0\Big\}
\end{equation}
where $\mathcal{T}$ is the set of bad triangles (i.e. triangles with exactly two positive edges and one negative edge). For $x \in \{0,1\}^{|E|}$, $x$ corresponds to disagreements in a clustering: we set $x_e = 1$ if $e$ is a disagreement and $x_e = 0$ otherwise. The constraints state that every clustering must make a disagreement on every bad triangle. Thus, ($P$) is a relaxation for the $\ell_1$-norm objective. In fact, we will prove the stronger statement that the fractional cost is $O(\textsf{OPT}_{P})$, where $\textsf{OPT}_{P}$ is the optimal objective value of ($P$). 

The dual is given by 
\begin{equation} \label{eq: dual_pivot} \tag{$D$}
\max\Big\{\sum_{T \in \mathcal{T}}y_T \mid \sum_{T \in \mathcal{T}: T \ni e} y_T \leq 1, \forall e \in E, y \geq 0 \Big\}.
\end{equation}

We show that by setting $y_T = \frac{1}{2\Delta}$ for all $T \in \mathcal{T}$, $y$ satisfies the following properties:
\begin{enumerate}
    \item $y$ is feasible in ($D$).
    \item The fractional cost of $d$ is at most $6 \cdot \sum_{T \in \mathcal{T}} y_T$.
\end{enumerate}

Letting $\textsf{OPT}_D$ be the optimal objective value of ($D$), we have $6 \cdot \sum_{T \in \mathcal{T}} y_T \leq 6 \cdot \textsf{OPT}_D = 6 \cdot \textsf{OPT}_{P} \leq 6 \cdot \textsf{OPT}$, which will conclude the proof.

To prove feasibility, we case on whether $e$ is positive or negative. 

If $e \in E^-$, then $|\{T \in \mathcal{T}: T \ni e\}| = |N_u^+ \cap N_v^+| \leq \Delta$, where equality is by the definition of a bad triangle. So 
\[\sum_{T \in \mathcal{T}: T \ni e} y_T  \leq \frac{\Delta}{2\Delta} \leq 1.\]


If $e \in E^+$, then $|\{T : T \ni e\}| = |(N_u^+ \cap N_v^-) \cup (N_u^- \cap N_v^+)| \leq 2\Delta$. We conclude $y$ is feasible, since
\[\sum_{T \in \mathcal{T}: T \ni e} y_T  \leq \frac{2\Delta}{2\Delta} = 1.\]
 

Now we need to show that the fractional cost of $d$ is bounded in terms of the objective value of $(D)$. First we bound the fractional cost of the negative edges:
\[\sum_{(u,v) \in E^-}(1-d_{uv}) \leq \sum_{(u,v)\in E^-} |N_u^+ \cap N_v^+|/\Delta= \sum_{e \in E^-} \sum_{T \in \mathcal{T}: T \ni e} 1/\Delta = \sum_{e \in E^-} \sum_{T \in \mathcal{T}: T \ni e} 2y_T,\]
where in the first inequality we have used that $|N_u^+ \cup N_v^+| \geq \Delta$. Next we bound the fractional cost of the positive edges:
\[\sum_{(u,v) \in E^+}d_{uv} \leq \sum_{(u,v)\in E^+} (|N_u^+ \cap N_v^-| + |N_u^- \cap N_v^+|)/\Delta = \sum_{e \in E^+} \sum_{T \in \mathcal{T}: T \ni e} 1/\Delta = \sum_{e \in E^+} \sum_{T \in \mathcal{T}: T \ni e} 2y_T. \]

So the total fractional cost is bounded by 
$\sum_{e \in E} \sum_{T \in \mathcal{T}: T \ni e} 2 y_T = 6 \cdot \sum_{T \in \mathcal{T}} y_T,$
since each triangle contains three edges. This is what we sought to show. Since the fractional cost of $d$ is bounded for the $\ell_1$-norm objective (and the $\ell_{\infty}$-norm objective by \cite{DMN23}), using $d$ as input to KMZ rounding algorithm produces a clustering that is simultaneously $O(1)$-approximate for the $\ell_1$- and $\ell_{\infty}$-norm objectives.  
\end{proof}


